\lecture{5.3}{Risk and Loss Functions}
\label{lec:5-3}

\subsection{Prediction Loss}

To compare regression estimators, we need a numerical score. A loss function measures
how bad a prediction is for one observation. If $\Qbar_n$ predicts the outcome using
$A$ and $W$, then a loss function assigns a number
\[
L(\Qbar_n)(W,A,Y)
\]
to the prediction-observation pair. Smaller loss means better prediction.

The choice of loss should match the regression target. If the target is the conditional
mean of a continuous outcome, squared-error loss is natural. If the target is the
conditional probability of a binary outcome, negative log-likelihood loss is natural.

\subsection{Squared-Error Loss}

For a continuous or mean-valued outcome regression, squared-error loss is
\[
L(\Qbar_n)(w,a,y)
=
\{y-\Qbar_n(a,w)\}^2.
\]
This loss heavily penalizes large prediction errors because the error is squared. It is
the standard loss for estimating conditional means.

\subsection{Negative Log-Likelihood Loss}

For a binary outcome $Y\in\{0,1\}$, an outcome regression estimates
\[
\Qbar(a,w)=\Prob(Y=1\given A=a,W=w).
\]
If $\Qbar_n(a,w)$ is a predicted probability, the Bernoulli likelihood contribution is
\[
\Qbar_n(a,w)^y\{1-\Qbar_n(a,w)\}^{1-y}.
\]
The negative log-likelihood loss is
\[
L(\Qbar_n)(w,a,y)
=
-\log\left[
\Qbar_n(a,w)^y\{1-\Qbar_n(a,w)\}^{1-y}
\right].
\]
Equivalently,
\[
L(\Qbar_n)(w,a,y)
=
-y\log\{\Qbar_n(a,w)\}
-(1-y)\log\{1-\Qbar_n(a,w)\}.
\]
This loss rewards assigning high probability to the outcome that actually occurred
and strongly penalizes confident but wrong predictions.

\subsection{Risk as Average Loss}

\begin{definition}[Risk]
The risk of a regression estimator $\Qbar_n$ under loss $L$ is the expected loss under
the true observed-data distribution:
\[
R(\Qbar_n)=\E\left\{L(\Qbar_n)(W,A,Y)\right\}.
\]
For discrete $(W,A,Y)$,
\[
R(\Qbar_n)
=
\sum_{w,a,y}
L(\Qbar_n)(w,a,y)\Prob(W=w,A=a,Y=y).
\]
\end{definition}

Risk is the average prediction error we would expect if the fitted regression were
evaluated on new observations from the same population. The best estimator is the
one with the smallest risk, provided the loss function is appropriate for the target.

\subsection{Valid Loss Functions}

\begin{definition}[Valid loss function]
A loss function is valid for a regression target $\Qbar$ if the true regression function
minimizes the population risk:
\[
\Qbar \in \argmin_Q R(Q).
\]
\end{definition}

This definition formalizes the idea that the scoring rule should reward estimators for
getting close to the true target. Squared-error loss is valid for conditional means.
Bernoulli negative log-likelihood loss is valid for conditional probabilities.

\subsection{Why the True Regression Minimizes Risk}

We now show why squared-error loss targets the conditional mean. Let
\[
\Qbar(A,W)=\E(Y\given A,W)
\]
be the true outcome regression, and let $Q(A,W)$ be any candidate prediction
function. The squared-error risk is
\[
R(Q)=\E\left[\{Y-Q(A,W)\}^2\right].
\]
Add and subtract the true regression:
\[
Y-Q(A,W)
=
\{Y-\Qbar(A,W)\}+\{\Qbar(A,W)-Q(A,W)\}.
\]
Squaring and taking expectations gives
\[
\begin{aligned}
R(Q)
&=
\E\left[
\{Y-\Qbar(A,W)\}^2
\right]\\
&\quad+
2\E\left[
\{Y-\Qbar(A,W)\}\{\Qbar(A,W)-Q(A,W)\}
\right]\\
&\quad+
\E\left[
\{\Qbar(A,W)-Q(A,W)\}^2
\right].
\end{aligned}
\]
The middle term is zero. To see this, use iterated expectation:
\[
\begin{aligned}
&\E\left[
\{Y-\Qbar(A,W)\}\{\Qbar(A,W)-Q(A,W)\}
\right]\\
&=
\E\left[
\{\Qbar(A,W)-Q(A,W)\}
\E\{Y-\Qbar(A,W)\given A,W\}
\right].
\end{aligned}
\]
But
\[
\E\{Y-\Qbar(A,W)\given A,W\}
=
\E(Y\given A,W)-\Qbar(A,W)=0.
\]
Therefore,
\[
R(Q)
=
\E\left[
\{Y-\Qbar(A,W)\}^2
\right]
+
\E\left[
\{\Qbar(A,W)-Q(A,W)\}^2
\right].
\]
The first term does not depend on $Q$. The second term is nonnegative and equals zero
only when $Q(A,W)=\Qbar(A,W)$ with probability one. Thus the true conditional mean
minimizes squared-error risk.

\begin{keyidea}
Risk is useful because an appropriate loss function turns regression estimation into a
well-defined competition: the estimator with lower risk is closer to the true
regression target in the sense encoded by the loss.
\end{keyidea}

\subsection{Risk Decomposition into Noise, Bias, and Variance}

Risk also clarifies the bias-variance tradeoff. Under squared-error loss,
\[
R(\Qbar_n)
=
\E\left[
\{Y-\Qbar(A,W)\}^2
\right]
+
\E\left[
\{\Qbar(A,W)-\Qbar_n(A,W)\}^2
\right].
\]
The first term is irreducible noise: even the true conditional mean cannot predict the
individual outcome perfectly.

Now decompose the second term. Add and subtract the repeated-sampling expectation
$\E\{\Qbar_n(A,W)\}$:
\[
\Qbar(A,W)-\Qbar_n(A,W)
=
\left[\Qbar(A,W)-\E\{\Qbar_n(A,W)\}\right]
+
\left[\E\{\Qbar_n(A,W)\}-\Qbar_n(A,W)\right].
\]
After squaring and taking expectations, the cross-term is zero by the definition of the
expectation of $\Qbar_n$. Thus
\[
\begin{aligned}
\E\left[
\{\Qbar(A,W)-\Qbar_n(A,W)\}^2
\right]
&=
\E\left[
\{\Qbar(A,W)-\E\{\Qbar_n(A,W)\}\}^2
\right] \\
&\quad+
\E\left[
\{\E\{\Qbar_n(A,W)\}-\Qbar_n(A,W)\}^2
\right].
\end{aligned}
\]
The first term is squared bias and the second is variance. Therefore squared-error
risk can be viewed as
\[
\text{risk}
=
\text{noise}+\text{bias}^2+\text{variance}.
\]

\begin{estimationbox}
Cross-validation will estimate risk using validation observations that were not used
to fit the candidate regression. This matters because empirical risk on the training
sample can reward overfitting and give an overly optimistic score.
\end{estimationbox}
